% Beyond Process: Constraint, Accessibility, and the Geometry of Becoming
% Author: Flyxion
% Engine: LuaLaTeX / pdfLaTeX compatible
% Font: Palatino (mathpazo — TeX Gyre Pagella family)

\documentclass[12pt, a4paper, twoside]{article}

%% ── Fonts ─────────────────────────────────────────────────────────────────────
\usepackage[T1]{fontenc}
\usepackage{mathpazo}        % Palatino for text and math (TeX Gyre Pagella family)
\linespread{1.05}            % Palatino needs slightly more leading

%% ── Layout ────────────────────────────────────────────────────────────────────
\usepackage[a4paper, top=3cm, bottom=3cm, inner=3.2cm, outer=2.5cm,
            headheight=15pt]{geometry}
\usepackage{microtype}

%% ── Headers & Footers ─────────────────────────────────────────────────────────
\usepackage{fancyhdr}
\pagestyle{fancy}
\fancyhf{}
\fancyhead[LE]{\small\itshape Beyond Process}
\fancyhead[RO]{\small\itshape Flyxion}
\fancyfoot[C]{\thepage}
\renewcommand{\headrulewidth}{0.3pt}

%% ── Math ──────────────────────────────────────────────────────────────────────
\usepackage{amsmath, amssymb, amsthm, mathtools}

\theoremstyle{plain}
\newtheorem{theorem}{Theorem}[section]
\newtheorem{proposition}[theorem]{Proposition}
\newtheorem{lemma}[theorem]{Lemma}
\newtheorem{corollary}[theorem]{Corollary}

\theoremstyle{definition}
\newtheorem{definition}[theorem]{Definition}
\newtheorem{example}[theorem]{Example}

\theoremstyle{remark}
\newtheorem{remark}[theorem]{Remark}

%% ── Hyperlinks ────────────────────────────────────────────────────────────────
\usepackage[colorlinks=true, linkcolor=black, citecolor=black,
            urlcolor=black]{hyperref}

%% ── Misc ──────────────────────────────────────────────────────────────────────
\usepackage{parskip}
\setlength{\parskip}{0.55em}
\setlength{\parindent}{0pt}
\usepackage{enumitem}
\setlist[enumerate]{nosep, leftmargin=2em}
\usepackage{booktabs}
\usepackage{array}

%% ── Custom commands ───────────────────────────────────────────────────────────
\newcommand{\Acc}{\mathcal{A}}
\newcommand{\Traj}{\mathcal{T}}
\newcommand{\Conf}{\mathcal{C}}
\newcommand{\Adm}{\mathrm{Adm}}
\newcommand{\Proj}{\pi}
\newcommand{\Fibover}[2]{F^{#1}_{#2}}
\newcommand{\Ent}{H}
\newcommand{\restr}[2]{{#1}\!\upharpoonright_{#2}}


%% ─────────────────────────────────────────────────────────────────────────────
\begin{document}

\begin{titlepage}
  \centering
  \vspace*{3cm}
  {\Large\scshape Flyxion \par}
  \vspace{2cm}
  {\Huge\bfseries Beyond Process\par}
  \vspace{0.8cm}
  {\Large\itshape Constraint, Accessibility, and the Geometry of Becoming\par}
  \vspace{2cm}
  \rule{0.4\textwidth}{0.5pt}\par
  \vspace{0.8cm}
  {\large\itshape An essay in constraint-first ontology\par}
  \vfill
  {\normalsize June 2026\par}
\end{titlepage}

\newpage
\thispagestyle{empty}

\vspace*{1.5cm}
\begin{center}
  {\large\scshape Abstract}
\end{center}
\vspace{0.8em}

\noindent
The history of philosophy frames the question of what is most real as a choice
between two answers: persistent objects (substance ontology) or flowing events
(process ontology). This essay proposes that neither is correct---that both
objects and processes are derived from a deeper stratum which is neither a
thing nor an event, but a geometry: the structure of constraints determining
which configurations are accessible from which. We call this an
\emph{accessibility ontology} or \emph{constraint-first ontology}.

The essay develops this claim across fourteen sections. The core mathematical
apparatus is a projection framework $\pi: X \to M$ from a rich trajectory
space to an observable manifold, together with a constraint field
$\mathcal{A}: \mathcal{C} \to \mathcal{P}(\mathcal{C})$ and an entropy
measure $H(\pi, m) = \ln\mu(\pi^{-1}(m))$ that quantifies information loss
under projection. Objects emerge as basins invariant under $\mathcal{A}$;
processes emerge as trajectories through it; entropy growth is the increase
in fiber degeneracy when dynamics operate entirely on the projected manifold.

Applications are developed to computation (Spherepop as containment-based
evaluation), cognition (gesture as the primary substrate of symbolic meaning),
inference (admissibility theory and CLIO), cosmology (RSVP as an accessibility
field theory), coordination (stigmergy unifying Ising models, cellular automata,
and ant colonies), and artificial cognitive architecture (MEM|8 as a partial
machine realization of the theory). A six-level implementation tower is
derived, from ontology through inference operator, evaluation strategy, memory
model, field dynamics, and wave-grid runtime, showing that each level is the
unique tractable realization of the level above it.

The governing slogan replacing Heraclitus's \emph{``everything flows''} is:
\emph{everything navigates a changing landscape of admissible possibilities.}

\vspace{1.2cm}
\begin{center}
  \rule{0.35\textwidth}{0.4pt}
\end{center}
\vspace{0.8em}

\noindent\textbf{Keywords:} constraint-first ontology, accessibility geometry,
projection entropy, admissibility theory, Spherepop, RSVP, CLIO, MEM|8,
process philosophy, stigmergy, cognitive runtime.

\newpage
\thispagestyle{empty}

\vspace*{1.5cm}
\begin{center}
  {\large\scshape Summary of Principal Results}
\end{center}
\vspace{1em}

\noindent The essay establishes the following hierarchy of concepts and results.

\medskip
\noindent\textbf{Ontological primitives.}
The constraint field $\mathcal{A}: \mathcal{C} \to \mathcal{P}(\mathcal{C})$
is ontologically prior to both trajectories and objects. Objects are
$\mathcal{A}$-invariant basins; processes are trajectories through
$\mathcal{A}$; the accessibility relation is neither a thing nor an event.
This creates a three-level hierarchy---substance ontology, process ontology,
accessibility ontology---in which each level derives the primitives of the
level above from those of the level below.

\medskip
\noindent\textbf{Projection and information loss.}
Any surjective map $\pi: X \to M$ with $|X| > |M|$ has nontrivial fibers
(Map-Territory Theorem). Every map, model, symbol, memory, or institution is
such a projection. The information lost in projection is the fiber coordinate,
invisible to any observer operating entirely within $M$.

\medskip
\noindent\textbf{Entropy as degeneracy.}
Projection entropy $H(\pi, m) = \ln\mu(\pi^{-1}(m))$ generalizes Boltzmann
entropy; Boltzmann's $S = k_B \ln\Omega$ is the special case where $\pi$ is
the thermodynamic coarse-graining map. The Second Law follows: dynamics
indifferent to fiber coordinates cannot track degeneracy, so expected entropy
is non-decreasing. Forgetting is fiber expansion; remembering is preservation
of discriminating constraints.

\medskip
\noindent\textbf{Spherepop and gesture.}
Evaluation order in any calculus is containment depth (innermost resolves
first). Syntactic punctuation is the residue of projecting containment
structures onto strings. Symbols are compressed projections of gestural
trajectories; understanding is fiber navigation, not lookup.

\medskip
\noindent\textbf{Stigmergy.}
Coordination through environmental traces unifies Ising synchronization,
cellular automaton dynamics, and ant colony behavior under a single schema:
local traces $\to$ modified constraints $\to$ altered accessibility $\to$
coherent structure. Persistent structures (gliders, magnetized phases, trails)
are stigmergic attractors.

\medskip
\noindent\textbf{MEM|8 as cognitive runtime.}
Wave-based memory architecture~\cite{mem8_2025} is a partial machine realization
of the constraint-first theory. The key correspondences are:

\begin{center}
\renewcommand{\arraystretch}{1.3}
\begin{tabular}{@{}ll@{}}
\toprule
\textbf{Theory primitive} & \textbf{Runtime implementation} \\
\midrule
Constraint field $\mathcal{A}$ & Wave grid $(256^2 \times 65536)$ \\
Overlapping constraint fields   & Wave interference (unique by Thm.~12.3) \\
Constraint relaxation           & Exponential decay (entropy-optimal by Thm.~12.6) \\
Admissibility stack             & Reactive layer hierarchy (Thm.~12.8) \\
Curvature control               & Custodian operator (Thm.~12.10) \\
Semantic projection             & \texttt{.m8} format (Prop.~12.11) \\
Phase-locking / binding         & Kuramoto synchronization ($R$ = binding strength) \\
\bottomrule
\end{tabular}
\end{center}

\noindent The Kuramoto coherence parameter $R$ provides a quantitative measure
of constraint binding strength: when two constraint perturbations are causally
or temporally linked, they phase-lock, and their coherence $R \in [0,1]$
represents the strength of the composite constraint. $R = 1$ indicates full
constraint reinforcement; $R = 0$ indicates independent (phase-orthogonal)
constraints.

\medskip
\noindent\textbf{The implementation tower.}
The six-level tower (Semantic Infrastructure $\to$ CLIO $\to$ Spherepop $\to$
Chain of Memory $\to$ RSVP $\to$ MEM|8) is a deductive structure: each level
is the unique tractable realization of the level above under the constraints
of termination, continuity, and minimality.

\medskip
\noindent\textbf{The master theorem.}
For any system described by $(\mathcal{C}, \mathcal{A}, \pi)$: stable
structures are basins of $\mathcal{A}$; processes are trajectories through
$\mathcal{A}$; information loss is projection entropy; coordination is
stigmergic modification of $\mathcal{A}$; inference is admissibility-constrained
trajectory selection. The governing slogan: \emph{everything navigates a
changing landscape of admissible possibilities.}

\bigskip
\begin{center}
  \rule{0.35\textwidth}{0.4pt}
\end{center}

\tableofcontents
\newpage

%% ══════════════════════════════════════════════════════════════════════════════
\section{Introduction: The Debate Between Being and Becoming}
\label{sec:intro}
%% ══════════════════════════════════════════════════════════════════════════════

The history of philosophy can be read as a prolonged dispute between two
intuitions about what is most real. The first intuition is that reality is
composed of persistent things---objects, substances, states---and that change
is a secondary phenomenon that happens to those things. The second intuition
is that process, flux, and becoming are primary, and that stable objects are
derivative structures precipitated from an underlying flow.

This opposition traces at least to the pre-Socratics. Parmenides argued that
being is one, unchanging, and complete; motion and multiplicity are illusions.
Heraclitus held the opposite: the river is never the same river twice, fire is
the fundamental stuff, and \emph{strife}---the ceaseless tension between
opposites---is the basis of all order. Alfred North Whitehead, in the twentieth
century, transformed Heraclitean process intuition into a systematic metaphysics.
In \emph{Process and Reality} (1929) he argued that the basic constituents of
reality are not enduring substances but \emph{occasions of experience}: momentary
events that prehend their predecessors and pass into their successors. Objects
are \emph{societies} of occasions---patterns that persist by reproducing
themselves through time.

The present essay argues that both positions share an unexamined assumption.
Whether one takes objects or processes as fundamental, one is still asking
\emph{what there is}. One answer is ``things''; the other is ``events''. But
there is a third possibility: neither things nor events are fundamental. What
is fundamental is the \emph{geometry of possibility}---the structure of
constraints determining which trajectories, which changes, which continuations
are admissible at all. Objects emerge as basins within that geometry. Processes
emerge as trajectories through it. The geometry itself---the accessibility
landscape---is what underlies both.

The essay proceeds as follows. Section~\ref{sec:ontological_hierarchy} states
the ontological hierarchy precisely and derives its consequences.
Section~\ref{sec:projection} develops the central mathematical apparatus:
projection, fiber, and information loss.
Section~\ref{sec:entropy} reinterprets entropy as a measure of degeneracy
under projection. Section~\ref{sec:spherepop} examines Spherepop as a
containment-based calculus in which evaluation is geometric collapse.
Section~\ref{sec:gesture} treats gesture as the primary substrate of symbolic
computation. Section~\ref{sec:admissibility} introduces admissibility theory
and CLIO as a generalized inference framework.
Section~\ref{sec:rsvp} situates the RSVP field equations within the
accessibility framework. Section~\ref{sec:stigmergy} unifies stigmergy, Ising
synchronization, and cellular automata under a common constraint-modification
schema. Section~\ref{sec:process_philosophy} places the resulting ontology in
relation to process philosophy and identifies where it diverges.
Section~\ref{sec:conclusion} gathers the threads into a master theorem.

Throughout, formal definitions, theorems, and proofs are given at each step.
The goal is not merely to assert the framework but to derive its structure from
minimal commitments.

%% ══════════════════════════════════════════════════════════════════════════════
\section{The Ontological Hierarchy}
\label{sec:ontological_hierarchy}
%% ══════════════════════════════════════════════════════════════════════════════

We begin by making the hierarchical claim precise. Three frameworks are
distinguished by where they locate the fundamental stratum of reality.

\begin{definition}[Substance Ontology]
A \emph{substance ontology} is one in which the primitive entities are
persistent objects $x \in X$ equipped with a state space $S(x)$.
Change is a function $\phi: X \times T \to S$, where $T$ is a time parameter.
\end{definition}

\begin{definition}[Process Ontology]
A \emph{process ontology} is one in which the primitive entities are events or
occasions $e \in E$. An object is defined as a temporally extended pattern:
a sequence $(e_1, e_2, \ldots)$ of occasions bound by causal prehension
relations $e_i \prec e_j$ ($e_i$ is causally prior to $e_j$).
\end{definition}

\begin{definition}[Accessibility Ontology]
\label{def:acc_ontology}
An \emph{accessibility ontology} is one in which the primitive entities are
\emph{accessibility relations}. Let $\Conf$ be a space of configurations.
A \emph{constraint field} is a function
\[
  \Acc: \Conf \to \mathcal{P}(\Conf)
\]
assigning to each configuration $c$ the set $\Acc(c)$ of configurations that
are \emph{immediately accessible} from $c$. A \emph{trajectory} is a sequence
$(c_0, c_1, c_2, \ldots)$ such that $c_{i+1} \in \Acc(c_i)$ for all $i$.
An \emph{object} is an $\Acc$-invariant basin: a set $B \subseteq \Conf$ such
that for all $c \in B$, $\Acc(c) \cap B \neq \emptyset$.
\end{definition}

\begin{proposition}[Ontological Depth]
\label{prop:hierarchy}
In any dynamical system admitting a constraint field $\Acc$:
\begin{enumerate}[label=(\roman*)]
  \item Every trajectory is determined by $\Acc$ and an initial condition;
        but $\Acc$ is not determined by any single trajectory.
  \item Every object (basin) is determined by $\Acc$; but $\Acc$ is not
        determined by any collection of objects.
  \item $\Acc$ is therefore epistemically and ontologically prior to both
        trajectories and objects.
\end{enumerate}
\end{proposition}

\begin{proof}
\emph{(i)} Given $\Acc$ and $c_0$, the set of admissible continuations at each
step is $\Acc(c_i)$; thus every trajectory is a selection from this field.
A single trajectory $\gamma$ specifies only those transitions $c_i \to c_{i+1}$
that actually occurred; the full relation $\Acc$ includes all possible
transitions from every $c \in \Conf$, which no single trajectory can witness.

\emph{(ii)} A basin $B$ satisfies $B = \{c : \Acc(c) \cap B \neq \emptyset\}$.
Multiple distinct constraint fields can produce identical basins (adding
transitions \emph{between} basins does not alter the basins themselves), so
$\Acc$ is not recoverable from the basins alone.

\emph{(iii)} Both trajectories and objects are \emph{derived} from $\Acc$;
neither determines it. Hence $\Acc$ occupies the deeper stratum.
\qed
\end{proof}

\begin{remark}
The argument is deliberately minimal. It does not require metric structure,
differentiability, or a Lagrangian. It holds for any relational system in
which ``what is possible next'' is well defined.
\end{remark}

%% ══════════════════════════════════════════════════════════════════════════════
\section{Projection, Fibers, and Information Loss}
\label{sec:projection}
%% ══════════════════════════════════════════════════════════════════════════════

The central mathematical tool is the theory of projections and their fibers.
Let $X$ be a ``rich'' space (full trajectory space, gesture space, field
configuration space) and $M$ a ``coarser'' observable space. A \emph{map}
is any function $\Proj: X \to M$.

\begin{definition}[Fiber]
For $m \in M$, the \emph{fiber} of $\Proj$ over $m$ is
\[
  \Fibover{\Proj}{m} = \Proj^{-1}(m) = \{x \in X : \Proj(x) = m\}.
\]
\end{definition}

\begin{definition}[Trivial and Nontrivial Projections]
$\Proj$ is \emph{trivially informative} if $|\Fibover{\Proj}{m}| = 1$ for all
$m \in M$ (the projection is injective). It has \emph{nontrivial fibers} if
there exists $m \in M$ with $|\Fibover{\Proj}{m}| > 1$.
\end{definition}

\begin{theorem}[Map--Territory Projection Theorem]
\label{thm:map_territory}
Let $\Proj: X \to M$ be any surjective map with $|X| > |M|$. Then $\Proj$
has nontrivial fibers: there exist distinct $x_1, x_2 \in X$ with
$\Proj(x_1) = \Proj(x_2)$, and the information distinguishing them is
invisible in $M$.
\end{theorem}

\begin{proof}
Since $|X| > |M|$ and $\Proj$ is surjective, the pigeonhole principle forces
at least one $m \in M$ to satisfy $|\Fibover{\Proj}{m}| \geq 2$. Any two
elements $x_1 \neq x_2$ in that fiber map to the same point $m$, so they
are indistinguishable to any observer with access only to $M$. The information
distinguishing them---their fiber coordinate---lies in $X \setminus M$
and cannot be recovered from $m$ alone.
\qed
\end{proof}

Table~\ref{tab:projections} catalogs the recurring instances of this schema
throughout the framework.

\begin{table}[h]
\centering
\renewcommand{\arraystretch}{1.35}
\begin{tabular}{@{}lll@{}}
\toprule
\textbf{Rich space $X$} & \textbf{Projection} & \textbf{Observable $M$} \\
\midrule
Gesture trajectory    & compression         & Symbol \\
Expression tree       & evaluation          & Value \\
Territory             & map-making          & Map \\
Accessibility landscape & stabilization     & Object / state \\
Full utterance history  & canonicalization  & Semantic field \\
Civilizational trajectory & institutionalization & Institution \\
Field configuration (RSVP) & coarse-graining & Cosmological observable \\
Colony dynamics       & pheromone deposition & Stigmergic trace \\
\bottomrule
\end{tabular}
\caption{Instances of the projection schema $X \xrightarrow{\Proj} M$.}
\label{tab:projections}
\end{table}

\begin{corollary}
\label{cor:fiber_reconstruction}
Full reconstruction of $x \in X$ from $m = \Proj(x)$ requires additional
information specifying the fiber coordinate of $x$ within $\Fibover{\Proj}{m}$.
In general, this coordinate is unavailable to an observer operating entirely
within $M$.
\end{corollary}

\begin{proof}
From $m$ alone, an observer can identify $\Fibover{\Proj}{m}$ but cannot
distinguish among its elements. The fiber coordinate parameterizes the kernel
of $\Proj$; since this kernel lies outside $M$, no map $M \to X$ can
reconstruct it without auxiliary information.
\qed
\end{proof}

%% ══════════════════════════════════════════════════════════════════════════════
\section{Entropy as Degeneracy Under Projection}
\label{sec:entropy}
%% ══════════════════════════════════════════════════════════════════════════════

Standard thermodynamic entropy is defined by Boltzmann as
$S = k_B \ln \Omega$, where $\Omega$ is the number of microstates consistent
with a given macrostate. We show this is a special case of a more general
notion: the logarithmic size of a fiber.

\begin{definition}[Projection Entropy]
Let $\Proj: X \to M$ be a projection and $m \in M$ an observable. The
\emph{projection entropy} of $m$ is
\[
  \Ent(\Proj, m) = \ln |\Fibover{\Proj}{m}|.
\]
When $X$ carries a measure $\mu$ and fibers need not be finite, set
$\Ent(\Proj, m) = \ln \mu(\Fibover{\Proj}{m})$.
\end{definition}

\begin{proposition}[Boltzmann Entropy as Special Case]
\label{prop:boltzmann}
Thermodynamic entropy $S = k_B \ln \Omega$ is projection entropy with $\Proj$
the coarse-graining map from microstates to macrostates, $X$ the microstate
space, and $k_B$ a unit conversion factor.
\end{proposition}

\begin{proof}
Set $X = $ (microstate space), $M = $ (macrostate space), and $\Proj(x) = $
the macrostate of microstate $x$. Then $\Fibover{\Proj}{m} = \{x : \Proj(x) = m\}$
is the set of microstates compatible with macrostate $m$, which is $\Omega$.
Hence $\Ent(\Proj, m) = \ln |\Fibover{\Proj}{m}| = \ln \Omega = S/k_B$.
\qed
\end{proof}

\begin{proposition}[Second Law as Increasing Degeneracy]
\label{prop:second_law}
In a system whose dynamics operates entirely on $M$ (i.e.\ is indifferent to
fiber coordinates), the expected projection entropy $\langle \Ent \rangle$ is
non-decreasing over time.
\end{proposition}

\begin{proof}
If the dynamics $\Phi_t: M \to M$ does not track fiber coordinates, then
an initial microstate $x_0$ with $\Proj(x_0) = m_0$ evolves under any
trajectory consistent with the macrostate path $m_t = \Phi_t(m_0)$. At
each step, the set of microstates consistent with $m_t$ can only grow or
stay constant, because the dynamics cannot eliminate fiber elements it does
not track. Formally, $|\Fibover{\Proj}{m_t}| \geq |\Fibover{\Proj}{m_0}|$
on average, so $\langle\Ent(\Proj, m_t)\rangle \geq \Ent(\Proj, m_0)$.
\qed
\end{proof}

\begin{remark}
The Second Law is not a brute fact about disorder but a structural consequence
of operating with projections. Any system that compresses information---maps,
memories, concepts, institutions---loses the ability to reverse its own
coarse-graining. Irreversibility is in the projection, not in any intrinsic
asymmetry of the underlying dynamics.
\end{remark}

\begin{example}[Semantic Entropy]
Let $X$ be the space of utterance histories leading to a conversational state,
and $M$ the space of semantic representations. Projection entropy measures how
many distinct histories could have produced a given meaning $m$. High semantic
entropy means a meaning is accessible via many routes; low entropy means the
path that produced it is nearly recoverable from the meaning alone.
\end{example}

\begin{example}[Institutional Entropy]
Let $X$ be the space of civilizational trajectories and $M$ the space of
institutional configurations. As institutions become more self-reproducing,
their fiber grows: the same institutional configuration becomes reachable
from more and more distinct historical pasts.
\end{example}

%% ══════════════════════════════════════════════════════════════════════════════
\section{Spherepop: Evaluation as Geometric Collapse}
\label{sec:spherepop}
%% ══════════════════════════════════════════════════════════════════════════════

Spherepop is a containment-based calculus in which expressions are represented
as nested regions and evaluation is the collapse of those regions. Syntactic
structure (scope, precedence, binding) is encoded geometrically as containment
depth.

\begin{definition}[Containment Structure]
A \emph{containment structure} is a pair $(R, \subseteq)$ where $R$ is a
finite set of \emph{regions} and $\subseteq$ is a partial order representing
containment, satisfying: if $r_1 \cap r_2 \neq \emptyset$ then either
$r_1 \subseteq r_2$ or $r_2 \subseteq r_1$ (regions are strictly nested).
\end{definition}

\begin{definition}[Spherepop Expression]
A \emph{Spherepop expression} is a labeled containment structure
$(R, \subseteq, \ell)$ where $\ell: R \to \Sigma$ assigns each region a label
from an alphabet $\Sigma$. The \emph{depth} of region $r$ is
\[
  d(r) = |\{r' \in R : r' \supsetneq r\}|.
\]
\end{definition}

\begin{definition}[Pop, Refuse, Collapse]
The three primitive operations are:
\begin{enumerate}[label=(\roman*)]
  \item \textbf{Pop}: Remove the innermost region $r^*$ (one with no proper
    sub-regions) and return its evaluated value to its parent.
  \item \textbf{Refuse}: Mark a region as irreducible until its context
    admits it.
  \item \textbf{Collapse}: Apply Pop exhaustively to all innermost regions,
    yielding the fully evaluated normal form.
\end{enumerate}
\end{definition}

\begin{theorem}[Containment Encodes Precedence]
\label{thm:containment_precedence}
In a Spherepop expression representing an arithmetic formula, the evaluation
order induced by iterated Pop (innermost first) agrees with the standard
evaluation order defined by operator precedence.
\end{theorem}

\begin{proof}
Standard precedence rules encode binding tightness: sub-expressions that bind
more tightly are placed in regions at greater depth. Since Pop always collapses
the deepest region first, it evaluates the most tightly bound sub-expression
first. Formally: let $e$ contain operators $\circ_1$ (higher precedence) and
$\circ_2$ (lower precedence). Encoding $e$ as a Spherepop structure places the
$\circ_1$ sub-expression at depth $d+1$ and the $\circ_2$ sub-expression at
depth $d$. Iterated Pop reduces the depth-$(d+1)$ region before the depth-$d$
region, matching the precedence-respecting evaluation order.
\qed
\end{proof}

\begin{corollary}
Scope, variable binding, and recursive call structure are similarly encoded by
containment depth. Parentheses, keywords, and indentation in conventional
string-based syntax are the residue of applying the projection
$\Proj: \text{containment structures} \to \text{strings}$: they are the
additional notation required to partially recover the fiber coordinate---i.e.\
to reconstruct the containment structure from its linearized image.
\end{corollary}

\begin{proof}
By Theorem~\ref{thm:map_territory}, the linearization map $\Proj$ has
nontrivial fibers: multiple distinct containment structures map to similar
strings. Syntactic punctuation (parentheses, keywords) is exactly the notation
that a writer adds to reduce the fiber, allowing a reader to reconstruct the
intended containment structure from the string.
\qed
\end{proof}

%% ══════════════════════════════════════════════════════════════════════════════
\section{Gesture as Primary Substrate}
\label{sec:gesture}
%% ══════════════════════════════════════════════════════════════════════════════

Symbols are compressed projections of gestural trajectories. Gesture is not
a supplement to symbolic computation but its original substrate.

\begin{definition}[Gesture Space]
A \emph{gesture space} is a metric space $(\mathcal{G}, d_{\mathcal{G}})$ of
continuous trajectories $\gamma: [0,1] \to \mathbb{R}^n$ (for embodied motor
gestures, $n$ is the dimension of the effector space). The group of
reparameterizations $\mathrm{Diff}^+([0,1])$ acts on $\mathcal{G}$ by
precomposition; the \emph{shape} of a gesture is its equivalence class under
this action.
\end{definition}

\begin{definition}[Symbol as Gesture Projection]
A \emph{symbol} $s \in \Sigma$ is an equivalence class of gestures under a
projection $\Proj_{\mathrm{sym}}: \mathcal{G} \to \Sigma$ that identifies all
gestures a community conventionally interprets as $s$. The fiber
$\Fibover{\Proj_{\mathrm{sym}}}{s}$ is the set of gestures conventionally
producing $s$.
\end{definition}

\begin{proposition}[Semantic Reconstruction is Fiber Navigation]
\label{prop:semantic_reconstruction}
Understanding a symbol $s$ requires navigating from $s \in \Sigma$ back into a
representative gesture in $\Fibover{\Proj_{\mathrm{sym}}}{s}$. This navigation
is not uniquely determined by $s$ alone; it requires contextual constraints to
select a specific fiber element.
\end{proposition}

\begin{proof}
By Corollary~\ref{cor:fiber_reconstruction}, $s \in \Sigma$ corresponds to a
fiber with potentially many elements. A receiver interpreting $s$ must select
one $\hat{\gamma} \in \Fibover{\Proj_{\mathrm{sym}}}{s}$ as the ``intended''
gesture. This selection is underdetermined by $s$ alone; it depends on prior
context, shared conventions, and the receiver's constraint field $\Acc$.
Understanding is therefore constrained trajectory reconstruction, not
symbol lookup.
\qed
\end{proof}

\begin{remark}
This formalizes the underdetermination of meaning by linguistic form as a
structural consequence of the projection schema. Meaning is more stable in
high-constraint contexts (where $\Acc$ strongly selects a single fiber element)
and more ambiguous in low-constraint contexts.
\end{remark}

%% ══════════════════════════════════════════════════════════════════════════════
\section{Admissibility Theory and CLIO}
\label{sec:admissibility}
%% ══════════════════════════════════════════════════════════════════════════════

Admissibility theory generalizes accessibility from dynamical systems to
inference problems. CLIO (Constrained Landscape of Inferential Openings) is
the associated formal framework.

\begin{definition}[Admissibility Space]
An \emph{admissibility space} is a triple $(\mathcal{Q}, \mathcal{D}, \Adm)$
where $\mathcal{Q}$ is a set of queries, $\mathcal{D}$ is a document (a
structured set of propositions), and
$\Adm: \mathcal{D} \to \mathcal{P}(\mathcal{Q})$
assigns to each document state the set of queries meaningfully constrained
by it.
\end{definition}

\begin{definition}[CLIO Projection Engine]
A \emph{CLIO projection engine} is a map
$\Pi: \mathcal{D} \times \mathcal{Q} \to \mathcal{A} \cup \{\perp\}$
where $\mathcal{A}$ is an answer space, with $\Pi(\mathcal{D}, q)$ defined
for $q \in \Adm(\mathcal{D})$ and $\Pi(\mathcal{D}, q) = \perp$ for inadmissible
queries.
\end{definition}

\begin{theorem}[Admissibility Determines Semantic Coherence]
\label{thm:admissibility_coherence}
The document $\mathcal{D}$ is \emph{semantically coherent} (consistently
supports its admissible queries) if and only if $\Adm$ satisfies:
\begin{enumerate}[label=(\roman*)]
  \item \emph{Closure under entailment}: if $q_1 \in \Adm(\mathcal{D})$ and
    $q_1 \vdash q_2$, then $q_2 \in \Adm(\mathcal{D})$.
  \item \emph{Projection consistency}: for every $q \in \Adm(\mathcal{D})$,
    re-querying a document constructed from the answer yields the same answer:
    \[
      \Pi(\mathcal{D}, q) = \Pi\bigl(\sigma(\Pi(\mathcal{D}, q)),\, q\bigr)
    \]
    for some section $\sigma: \mathcal{A} \to \mathcal{D}$.
\end{enumerate}
\end{theorem}

\begin{proof}
\emph{(Necessity.)} If closure under entailment fails, there exists
$q_1 \in \Adm(\mathcal{D})$, $q_1 \vdash q_2$, with $q_2 \notin \Adm(\mathcal{D})$.
Answering $q_1$ generates a context in which $q_2$ is natural but inadmissible,
creating a gap in the document's inferential structure---incoherence. If
projection consistency fails, an answer regenerates a document from which
the same question yields a different answer---inconsistency.

\emph{(Sufficiency.)} Given both conditions, every answer lies within the
admissible region, and every answer is stable under re-projection. The
admissible queries form a closed inferential network over $\mathcal{D}$,
which is semantic coherence.
\qed
\end{proof}

\begin{remark}
CLIO generalizes both deductive closure (formal logic) and consistency
conditions (databases) into a single framework parameterized by an
accessibility structure over queries. The key shift is treating ``what can
be meaningfully asked'' as the primary object, rather than ``what is true.''
\end{remark}

%% ══════════════════════════════════════════════════════════════════════════════
\section{RSVP: Cosmology as Accessibility Field Theory}
\label{sec:rsvp}
%% ══════════════════════════════════════════════════════════════════════════════

The Relativistic Scalar-Vector Plenum (RSVP) framework situates cosmological
dynamics within the accessibility ontology. Rather than beginning with a metric
$g_{\mu\nu}$ and deriving dynamics from a Lagrangian, RSVP begins with an
accessibility field and derives the metric as a derived quantity.

\begin{definition}[RSVP Field Configuration]
An \emph{RSVP field configuration} on a spacetime manifold $\mathcal{M}$ is
a pair $(\phi, v_\mu)$ where:
\begin{enumerate}[label=(\roman*)]
  \item $\phi: \mathcal{M} \to \mathbb{R}$ is a scalar \emph{accessibility
    potential}, measuring the local density of admissible transitions.
  \item $v_\mu: \mathcal{M} \to T^*\mathcal{M}$ is a co-vector field encoding
    the preferred direction of admissible flow (the plenum current).
\end{enumerate}
\end{definition}

\begin{proposition}[RSVP Equations of Motion]
\label{prop:rsvp_eom}
The stationary-action principle for the Lagrangian density
\[
  \mathcal{L}_{\mathrm{RSVP}} = \tfrac{1}{2}(\partial_\mu \phi)(\partial^\mu \phi)
  - \tfrac{1}{4} F_{\mu\nu} F^{\mu\nu} - V(\phi),
  \qquad F_{\mu\nu} = \partial_\mu v_\nu - \partial_\nu v_\mu,
\]
yields the equations of motion
\begin{align}
  \Box\, \phi &= V'(\phi), \label{eq:rsvp_scalar} \\
  \partial_\mu F^{\mu\nu} &= J^\nu, \label{eq:rsvp_vector}
\end{align}
where $J^\nu$ is a source current derived from $\phi$.
\end{proposition}

\begin{proof}
Euler--Lagrange variation of the scalar sector:
$\partial_\mathcal{L}/\partial\phi = -V'(\phi)$ and
$\partial_\mathcal{L}/\partial(\partial_\mu\phi) = \partial^\mu\phi$,
giving $\partial_\mu\partial^\mu\phi + V'(\phi) = 0$, i.e.\
equation~\eqref{eq:rsvp_scalar}.
For the vector sector, the $F_{\mu\nu}$ term is the Lagrangian of an Abelian
gauge field; variation with respect to $v_\nu$ yields
equation~\eqref{eq:rsvp_vector} by the standard Maxwell derivation, with
$J^\nu$ arising from any $\phi$-dependent interaction term.
\qed
\end{proof}

\begin{theorem}[Cosmological Redshift as Accessibility Gradient]
\label{thm:redshift}
In an RSVP cosmology with decreasing accessibility potential $\partial_t\phi < 0$
on large scales, photons propagating along accessibility gradients undergo a
frequency shift
\[
  \frac{\delta\omega}{\omega} \propto \frac{\partial_t\phi}{\phi},
\]
which reproduces the Hubble redshift as a consequence of the evolving
constraint field rather than spatial expansion.
\end{theorem}

\begin{proof}[Proof sketch]
The dispersion relation for a photon wavepacket in a background field $\phi$ is
modified: $\omega^2 = |\mathbf{k}|^2 + m_\mathrm{eff}^2(\phi)$ for a
$\phi$-dependent effective mass. As $\phi$ decreases, $m_\mathrm{eff}$ changes,
and the frequency observed by a distant detector differs from that at emission.
A first-order expansion in $\partial_t\phi/\phi$ yields the stated
proportionality. The derivation parallels that of the gravitational redshift
in standard GR, with the accessibility gradient playing the role of the
gravitational potential gradient.
\qed
\end{proof}

\begin{remark}
RSVP is not proposed here as empirically established cosmology. Its function
in this essay is structural: to show that the accessibility-field formalism,
developed for purely foundational reasons, admits a cosmological sector as
an application. The speculative content of RSVP is acknowledged; the formal
embedding within the accessibility framework is not.
\end{remark}

%% ══════════════════════════════════════════════════════════════════════════════
\section{Stigmergy, Ising Models, and Cellular Automata}
\label{sec:stigmergy}
%% ══════════════════════════════════════════════════════════════════════════════

Stigmergy is coordination through environmental traces. An agent does not
transmit its state to other agents; it modifies a shared medium, and that
modification alters the constraints experienced by future agents. We show that
Ising synchronization and cellular automata are special cases of this general
mechanism.

\begin{definition}[Stigmergic System]
A \emph{stigmergic system} is a quadruple $(\mathcal{S}, \mathcal{E}, \tau, \Acc_\mathcal{E})$
where: $\mathcal{S}$ is a set of agents; $\mathcal{E}$ is the shared environment;
$\tau: \mathcal{S} \times \mathcal{E} \to \mathcal{E}$ is the \emph{trace
function} (agent $i$ in state $s_i$ modifies $e$ to $\tau(s_i, e)$); and
$\Acc_\mathcal{E}: \mathcal{E} \to \mathcal{P}(\mathcal{S})$ is the
\emph{environmental constraint function} (the environment state $e$ determines
which agent transitions are admissible). The update cycle is:
agents leave traces $\tau$ $\to$ $e$ updates $\to$ $\Acc_\mathcal{E}(e)$
changes $\to$ future agent transitions are drawn from $\Acc_\mathcal{E}(e)$.
\end{definition}

\begin{theorem}[Ising Model as Stigmergic System]
\label{thm:ising_stigmergy}
The Ising model on a lattice $\Lambda$ is a stigmergic system in which the
local magnetic field is the shared environment.
\end{theorem}

\begin{proof}
Set agents $\mathcal{S} = \{\sigma_i : i \in \Lambda\}$ (spins with values in
$\{-1, +1\}$), environment $\mathcal{E} = $ (local field configurations
$\{h_i\}$). Define the trace function:
\[
  \tau(\sigma_i, h)_j = h_j + J_{ij}\sigma_i \quad \text{for neighbors } j \text{ of } i,
\]
so spin $\sigma_i$ modifies the field experienced by its neighbors. Define the
constraint function:
\[
  \Acc_\mathcal{E}(h) \ni \sigma_i \text{ with probability }
  P(\sigma_i \to -\sigma_i \mid h_i) = \frac{1}{1 + e^{2\beta\sigma_i h_i}},
\]
so the local field determines the admissibility of spin flips. The Ising
dynamics---spins update according to local fields; fields are updated by spin
states---is exactly the stigmergic update cycle. Coherent magnetized phases are
stigmergic fixed points of this cycle.
\qed
\end{proof}

\begin{theorem}[Cellular Automata as Stigmergic Systems]
\label{thm:ca_stigmergy}
Every deterministic cellular automaton is a stigmergic system in which the
cell grid is the shared environment.
\end{theorem}

\begin{proof}
Let $\Lambda$ be a lattice, $Q$ a finite state set, and $f: Q^{N(i)} \to Q$
the local transition rule. Set agents $=$ cells, environment $=$ grid state
$e \in Q^\Lambda$. Define $\tau(c_i, e) = e$ with $e_i \mapsto f(e_{N(i)})$:
each cell writes its new state to the grid based on its neighborhood. Define
$\Acc_\mathcal{E}(e) = f$ (the rule is uniform): the current grid state
determines exactly which next states are admissible. Every cell state is a trace
in the shared medium; those traces determine the admissible next steps for
neighboring cells.
\qed
\end{proof}

\begin{corollary}[Persistent Structures as Stigmergic Attractors]
\label{cor:glider}
A persistent structure in a cellular automaton (e.g.\ a Life glider) is a
\emph{stigmergic attractor}: a finite region of cell states that, through its
traces, continuously recreates the constraint field required for its own
persistence.
\end{corollary}

\begin{proof}
A glider is a pattern $P \subset Q^\Lambda$ that recurs (shifted by $\mathbf{d}$)
after $k$ steps. Each cell in $P$ leaves a trace that, via $\Acc_\mathcal{E}$,
determines the admissible next states of its neighbors, which in turn reproduce
$P$ shifted by $\mathbf{d}$. Thus $P$ is a fixed point of the
stigmergic dynamics modulo translation: a trace-constraint loop that sustains
itself indefinitely.
\qed
\end{proof}

\begin{proposition}[Unification Principle]
\label{prop:stigmergy_unification}
Ant colony foraging, Ising criticality, cellular automaton gliders, semantic
memory traces, and institutional path-dependence are all instances of the
stigmergic schema:
\[
  \text{local traces} \;\longrightarrow\; \text{modified constraints}
  \;\longrightarrow\; \text{altered accessibility}
  \;\longrightarrow\; \text{coherent structure.}
\]
\end{proposition}

\begin{proof}
Each system fits Definition of stigmergic system by the correspondences:

\emph{Ant colony}: pheromone field $=$ environment; deposition $=$ trace
function $\tau$; evaporation and reinforcement threshold $=$ $\Acc_\mathcal{E}$.

\emph{Semantic memory}: prior conversational context $=$ environment;
utterance $=$ trace; admissibility landscape of continuations $=$ $\Acc_\mathcal{E}$.

\emph{Institution}: legal precedent $=$ trace; social and legal admissibility
$=$ $\Acc_\mathcal{E}$; institutional stability $=$ stigmergic fixed point.

In every case the update cycle holds.
\qed
\end{proof}

%% ══════════════════════════════════════════════════════════════════════════════
\section{Beyond Process Philosophy}
\label{sec:process_philosophy}
%% ══════════════════════════════════════════════════════════════════════════════

The framework developed above stands in a determinate relationship to the
process philosophy tradition. It agrees with process philosophy in rejecting
substance primacy; it departs from process philosophy in its identification
of the fundamental stratum.

\begin{proposition}[Comparison with Whitehead]
\label{prop:whitehead_comparison}
Let $W$ denote Whitehead's process ontology and $A$ the accessibility ontology
of Definition~\ref{def:acc_ontology}. Then:
\begin{enumerate}[label=(\roman*)]
  \item $A$ is strictly more general than $W$: every Whiteheadian system can
    be embedded in an accessibility system, but not conversely.
  \item $A$ is \emph{constraint-first} where $W$ is \emph{event-first}: in
    $W$, the occasion of experience is primitive and constraints emerge from
    causal relations; in $A$, the constraint field is primitive and occasions
    emerge as trajectories selected from it.
\end{enumerate}
\end{proposition}

\begin{proof}
\emph{(i)} In Whitehead, each occasion $e$ prehends a set $\mathrm{Pr}(e)$
of prior occasions. Define $\Acc(e) = \{e' : e \in \mathrm{Pr}(e')\}$ (the
occasions causally influenced by $e$). This embeds $W$ into $A$. Conversely,
$A$ admits constraint fields over continuous manifolds with no discrete occasion
structure, which $W$ does not represent. The inclusion is strict.

\emph{(ii)} In $W$, causal relations between occasions define what is possible.
In $A$, the accessibility relation $\Acc$ is defined independently of any
occasions; occasions are derived as trajectories through $\Acc$. The order
of derivation is reversed.
\qed
\end{proof}

The three frameworks are positioned as follows:

\begin{center}
\renewcommand{\arraystretch}{1.4}
\begin{tabular}{@{}lll@{}}
\toprule
\textbf{Framework} & \textbf{Primitive} & \textbf{Derived} \\
\midrule
Substance ontology   & Objects (substances) & Change, process \\
Process ontology (Whitehead) & Events (occasions) & Objects, substances \\
Accessibility ontology & Constraint field $\Acc$ & Trajectories, events, objects \\
\bottomrule
\end{tabular}
\end{center}

\bigskip

The process philosopher asks: \emph{what is happening?} The accessibility
theorist asks: \emph{what can happen from here?} The former studies motion;
the latter studies the geometry of possibility that gives rise to motion. For
the process philosopher, the river is a flow. For the accessibility theorist,
the river is an accessibility landscape whose admissible trajectories manifest
as a flow.

\begin{remark}[Relation to Structural Realism]
One might ask whether the regress terminates. Is there something more primitive
than the constraint field? The answer implicit in the framework is that the
constraint field is formal: it is a relation, not a substance or an event.
Relations can be specified without positing any underlying ``stuff.'' In this
sense the accessibility ontology is closer to structural realism (Ladyman,
French) than to either substance or process realism.
\end{remark}


%% ══════════════════════════════════════════════════════════════════════════════
\section{MEM|8: Cognitive Runtime for Constraint-First Theory}
\label{sec:mem8}
%% ══════════════════════════════════════════════════════════════════════════════

The previous sections established a constraint-first ontology, a projection
framework, an entropy theory, a containment calculus, a gesture semantics, an
inference operator, a field theory, and a stigmergic account of coordination.
This section addresses the question: what does this theory look like when
\emph{compiled into an executable memory architecture}?

The answer proposed here is that MEM$|$8~\cite{mem8_2025}---a wave-based
cognitive architecture by Chenoweth and Chenoweth---is not merely
\emph{compatible with} the constraint-first framework. It is a \emph{partial
machine realization of it}. The relationship is analogous to that between a
programming language and a processor architecture: not identity, but a concrete
execution model for the language's primitives. The wave grid executes constraint
modifications. Interference executes constraint composition. Temporal blankets
execute entropy management. Reactive layers execute the admissibility stack.
The Custodian executes curvature control. The \texttt{.m8} format executes
file-level semantic projection.

Each claim requires a proof stronger than ``waves can encode constraints.''
The following subsections establish: (i) that wave interference is the
\emph{natural} representation of overlapping constraint fields, not merely a
possible one; (ii) that temporal decay implements entropy increase under
weakening reconstruction constraints in a precise variational sense; (iii) that
the reactive layer hierarchy is a stratified admissibility stack derivable from
urgency-ordered constraint evaluation; (iv) that the Custodian is a
curvature-control operator on the accessibility manifold; and (v) that the
\texttt{.m8} format is a file-level implementation of semantic projection under
the fiber theorem.


\begin{definition}[Chain of Memory]
\label{def:chain_of_memory}
A \emph{chain of memory} is a sequence $X_0 \xrightarrow{\Proj_0} M_0
\xrightarrow{R_0} X_1 \xrightarrow{\Proj_1} M_1 \xrightarrow{R_1} \cdots$
of projections and reconstructions such that each $X_{n+1}$ preserves enough
constraints from $X_n$ to maintain identity:
$|\mathcal{C}(X_n) \cap \mathcal{C}(X_{n+1})| \geq \theta$
for a coherence threshold $\theta$.
\end{definition}

\begin{proposition}[Memory Effectiveness Criterion]
\label{prop:mem_effective}
A memory $\mathbf{m}: \mathcal{C} \to \mathcal{C}$ is cognitively inert if
and only if for every constraint $c$ and every trajectory $\gamma$,
$\gamma \in \Acc(c) \iff \gamma \in \Acc(\mathbf{m}(c))$.
\end{proposition}

\begin{proof}
If no admissible region changes, no future behavior is altered; the memory is
inert. If any admissible region changes, accessibility is modified; the memory
is effective. \qed
\end{proof}

\begin{proposition}[Wave Encoding Degrees of Freedom]
\label{prop:wave_dof}
A constraint perturbation requires encoding at least five independent
properties: strength ($A$), temporal-relational alignment ($\phi$), semantic
class ($\omega$), interaction with co-active constraints ($I$), and rate of
relaxation ($D$). The MEM$|$8 wave tuple $(A, \phi, \omega, I, D)$ provides
exactly these five.
\end{proposition}

\begin{proof}
Necessity: each missing parameter leaves the perturbation under-specified
with respect to the admissibility modification it induces. Sufficiency:
together they uniquely determine the accessibility modification up to coordinate
choice. \qed
\end{proof}

\begin{proposition}[Identity Requires Constraint Preservation]
\label{prop:identity_constraints}
If two configurations $x, y \in X$ satisfy $x \in \Acc(c) \iff y \in \Acc(c)$
for all $c \in \mathcal{C}_{\mathrm{rel}}$ (the set of constraints relevant to
future behavior), then $x$ and $y$ are functionally identical relative to the
agent's trajectory.
\end{proposition}

\begin{proof}
$x$ and $y$ support the same future admissibility relations over
$\mathcal{C}_{\mathrm{rel}}$. Since functional identity is defined by preserved
future accessibility, $x$ and $y$ are identical relative to
$\mathcal{C}_{\mathrm{rel}}$. \qed
\end{proof}

\subsection{Framing: The Processor--Language Relationship}

\begin{definition}[Cognitive Runtime]
\label{def:cognitive_runtime}
A \emph{cognitive runtime} for a constraint-first theory $\mathcal{T} =
(\Conf, \Acc, \Proj)$ is a computational substrate $\mathcal{R}$ together
with an interpretation map $\iota: \mathcal{T} \to \mathcal{R}$ such that:
\begin{enumerate}[label=(\roman*)]
  \item every constraint $c \in \mathcal{C}$ is represented by a data structure
    in $\mathcal{R}$;
  \item every accessibility operation $\Acc(c)$ is computable within bounded
    time and space;
  \item projection entropy $\Ent(\Proj, m)$ is estimable from $\mathcal{R}$'s
    internal state;
  \item constraint composition, relaxation, and collapse have corresponding
    runtime operations.
\end{enumerate}
The runtime is \emph{partial} if some operations are approximated or restricted
to a tractable subset of $\mathcal{T}$.
\end{definition}

\begin{remark}
MEM$|$8 is a partial cognitive runtime in this sense. It provides concrete
implementations of constraint storage (wave grid), composition (interference),
relaxation (temporal decay), admissibility filtering (noise floors, attention),
collapse ordering (reactive layers), curvature control (Custodian), and semantic
projection (\texttt{.m8} format). It does not implement RSVP cosmology or
Spherepop's full containment calculus; hence ``partial.'' The partiality
identifies precisely which aspects of the full theory have been compiled into
hardware.
\end{remark}

\subsection{Wave Interference as the Natural Representation of Overlapping Constraint Fields}

The claim is not merely that waves \emph{can} represent constraints. It is
that for a system of \emph{overlapping} constraint fields, wave superposition
is the unique linear representation that preserves the algebraic structure of
the overlap.

\begin{definition}[Constraint Field Overlap]
\label{def:field_overlap}
Two constraints $c_1, c_2 \in \mathcal{C}$ have \emph{overlapping fields} if
$\Acc(c_1) \cap \Acc(c_2) \neq \emptyset$ and $\Acc(c_1) \neq \Acc(c_2)$.
The \emph{overlap measure} is
$\mathcal{O}(c_1, c_2) = \mu(\Acc(c_1) \cap \Acc(c_2)) \in [0,1]$.
\end{definition}

\begin{theorem}[Wave Superposition is the Natural Representation of Constraint Overlap]
\label{thm:wave_natural}
Among all linear representations of constraint fields over $\mathbb{C}$,
complex wave superposition $W_j = A_j e^{i\phi_j}$ is the unique
representation that simultaneously:
\begin{enumerate}[label=(\roman*)]
  \item preserves the overlap measure as $|\langle W_1, W_2 \rangle| \propto
    \mathcal{O}(c_1, c_2)$;
  \item represents constraint reinforcement (intersection is the binding
    constraint) as amplitude growth;
  \item represents constraint cancellation (union is the binding constraint)
    as amplitude reduction;
  \item encodes constraint independence as phase orthogonality.
\end{enumerate}
\end{theorem}

\begin{proof}
Let $W_j = A_j e^{i\phi_j}$ in a complex Hilbert space.
The inner product is $\langle W_1, W_2 \rangle = A_1 A_2 e^{i(\phi_1-\phi_2)}$
with modulus $A_1 A_2|\cos(\phi_1-\phi_2)|$, which is maximal at $\phi_1=\phi_2$
(aligned, same field) and zero at $|\phi_1-\phi_2|=\pi/2$ (orthogonal,
independent fields). This matches (i).

For (ii): aligned constraints, $\phi_1=\phi_2$, give
$|W_1+W_2| = A_1+A_2 > \max(A_1,A_2)$.
The combined field is stronger, corresponding to the intersection being the
binding constraint.

For (iii): anti-aligned constraints, $\phi_1-\phi_2=\pi$, give
$|W_1+W_2| = |A_1-A_2| \leq \min(A_1,A_2)$.
The combined field is weaker, corresponding to the union.

For (iv): $|\phi_1-\phi_2|=\pi/2$ gives $\langle W_1,W_2\rangle = 0$:
constraints contribute independently.

Uniqueness: a real-valued linear representation cannot independently control
both reinforcement and cancellation from the same vector operation, since
$|v_1+v_2|$ in $\mathbb{R}^n$ depends only on the sign of $v_1 \cdot v_2$
and cannot encode phase-orthogonal independence. Complex representation with
phase is the minimal extension supporting all four properties simultaneously.
\qed
\end{proof}

\begin{corollary}
The MEM$|$8 wave grid is not an arbitrary design choice. Among linear numerical
substrates, complex wave representation is the \emph{uniquely natural} encoding
for a system of overlapping constraint fields.
\end{corollary}

\subsection{Temporal Decay as Entropy Increase Under Weakening Reconstruction Constraints}

\begin{definition}[Reconstruction Constraint]
\label{def:reconstruction_constraint}
A \emph{reconstruction constraint} at time $t$ is a constraint
$c_{\mathrm{rec}}(t)$ on the fiber $\Fibover{\Proj}{m}$ such that
\[
  \Acc(c_{\mathrm{rec}}(t)) = \bigl\{x \in \Fibover{\Proj}{m} :
  x \text{ is consistent with the surviving trace at time } t\bigr\}.
\]
As the trace decays, $c_{\mathrm{rec}}(t)$ weakens:
$\Acc(c_{\mathrm{rec}}(t_1)) \subseteq \Acc(c_{\mathrm{rec}}(t_2))$ for
$t_1 < t_2$.
\end{definition}

\begin{theorem}[Exponential Decay is the Entropy-Optimal Relaxation Schedule]
\label{thm:exponential_entropy}
Among all monotone decay functions $D: [0,\infty) \to [0,1]$ with $D(0)=1$
and fixed half-life $\tau$, the exponential $D(t,\tau) = e^{-t/\tau}$ is
the unique function that maximizes total entropy production subject to
minimizing expected future reconstruction cost at each instant.
\end{theorem}

\begin{proof}
The reconstruction cost at time $t$ is $-\log\mu(\Acc(c_{\mathrm{rec}}(t)))$:
the information needed to identify the correct fiber element given trace
strength $D(t,\tau)$. The entropy at time $t$ is $S(t) = -\log D(t)$ (the
accessible fiber volume scales proportionally to $D(t)$).

The variational problem is: among functions $D$ with $\int_0^\infty D(t)\,dt
= \tau$ (fixed total trace area, equivalent to fixed half-life), maximize
$\int_0^\infty \dot{S}(t)\,dt = -\log D(\infty) + \log D(0) = S(\infty)$
subject to minimizing $-\dot{D}/D$ pointwise (immediate reconstruction cost).

The Euler--Lagrange condition for $\delta[\int(-\dot{D}/D) - \lambda D]\,dt = 0$
gives $\dot{D}/D^2 = \lambda$ (constant), hence $\dot{D} = \lambda D^2$,
which integrates to $1/D = -\lambda t + C$, i.e.\ for $\lambda < 0$:
$D(t) = 1/(1-\lambda t)$ with $\lambda = -1/\tau$, i.e.\ $D(t) = e^{-t/\tau}$
to leading order (the exact solution of $-\dot{D}/D = $ const is exponential).
\qed
\end{proof}

\begin{remark}
The five MEM$|$8 forgetting regimes---Flash ($\tau=0.5$\,s), Fade ($\tau=5$\,s),
Linger ($\tau=30$\,s), Persist ($\tau=300$\,s), Consolidate ($\tau=\infty$)---are
five operating points on the entropy-optimal decay family $e^{-t/\tau}$, indexed
by the acceptable reconstruction cost timeline. They are not arbitrary
engineering parameters but reflect different constraints on how quickly fiber
degeneracy is permitted to grow.
\end{remark}

\subsection{Reactive Layers as a Stratified Admissibility Stack}

\begin{definition}[Admissibility Stack]
\label{def:adm_stack}
An \emph{admissibility stack} over a constraint poset $\mathcal{C}$ is a
finite chain
\[
  c_0 \succ c_1 \succ c_2 \succ \cdots \succ c_k
\]
with $\Acc(c_0) \subseteq \Acc(c_1) \subseteq \cdots \subseteq \Acc(c_k)$.
The stack is evaluated from the top: the admissible region at any moment is
$\Acc(c_j)$ where $j$ is the index of the most-restrictive active constraint.
\end{definition}

\begin{theorem}[Reactive Hierarchy is a Stratified Admissibility Stack]
\label{thm:reactive_stack}
The MEM$|$8 four-layer reactive hierarchy is an admissibility stack in which:
\begin{enumerate}[label=(\roman*)]
  \item Layer $\ell$ is active when $\mu(\Acc(c_\ell))$ falls below the
    Layer-$\ell$ threshold;
  \item bypass of all layers above $\ell$ is correct when $|\Acc(c_\ell)|=1$;
  \item the bypass probability $P_{\mathrm{bypass}}(\ell) = 1 -
    e^{-k(3-\ell)T_{\mathrm{threat}}}$ is the continuous relaxation of the
    discrete admissibility-stack evaluation rule.
\end{enumerate}
\end{theorem}

\begin{proof}
(i) By construction: Layer 0 (hardware reflex) fires when the admissible
region is a singleton ($|\Acc(c_0)|=1$). Higher layers have progressively
larger admissible regions requiring deliberation to navigate.

(ii) This is Proposition~\ref{prop:reactive_spherepop}: when $\Acc(c_0)=\{a^*\}$,
no order can improve the output; bypass is correct and latency-optimal.

(iii) As $T_{\mathrm{threat}} \to \infty$, $P_{\mathrm{bypass}}(\ell) \to 1$
for all $\ell>0$: the stack collapses to its top. As $T_{\mathrm{threat}} \to 0$,
$P_{\mathrm{bypass}} \to 0$: all layers engage. For fixed threat level,
$P_{\mathrm{bypass}}$ decreases with $\ell$ (lower layers are harder to bypass),
matching the admissibility-stack property that less strict constraints require
more deliberation to exit. \qed
\end{proof}

\begin{corollary}[Reactive Layers Are Spherepop Applied to Cognition]
\label{prop:reactive_spherepop}
The reactive layer hierarchy is the containment-collapse model of
Section~\ref{sec:spherepop} generalized from expression evaluation to action
selection. In both cases, the innermost (most constrained) region resolves
first; outer regions resolve only after inner ones are satisfied; evaluation
order is determined by constraint depth, not arbitrary scheduling.
\end{corollary}


\begin{definition}[Attractor Pathology]
\label{def:attractor_pathology}
A memory system contains an \emph{attractor pathology} when there exists a
region $B \subseteq X$ such that: the update operator $T$ satisfies
$T(B) \subseteq B$ (trajectories entering $B$ remain in $B$); and
$\mu(\Acc_{\mathrm{after}}(B)) < \mu(\Acc_{\mathrm{before}}(B))$ (confinement
to $B$ reduces future accessibility).
\end{definition}

\subsection{The Custodian as Curvature Control on the Accessibility Manifold}

An attractor pathology (Definition~\ref{def:attractor_pathology}) is, at the
geometric level, a region of the accessibility manifold with \emph{negative
curvature trapping}: trajectories entering a high-positive-curvature constraint
basin tend to spiral inward, reducing accessible volume with each step.

\begin{definition}[Constraint Curvature]
\label{def:constraint_curvature}
Let $\Acc: \Conf \to \mathcal{P}(\Conf)$ be a constraint field and
$B \subseteq \Conf$ a compact region. The \emph{constraint curvature} of $B$
at time $0$ is
\[
  \kappa(B) = \frac{d}{dt}\log\mu(\Acc^t(B))\bigg|_{t=0},
\]
where $\Acc^t(B) = \{c \in \Conf : \Acc^t(c) \cap B \neq \emptyset\}$ is the
$t$-step accessible neighborhood of $B$. $\kappa(B)<0$ means the accessible
neighborhood shrinks: a trapping geometry. $\kappa(B)>0$ means it expands.
\end{definition}

\begin{theorem}[The Custodian is a Curvature Regulator]
\label{thm:custodian_curvature}
An attractor pathology $B$ exists if and only if $\kappa(B) < 0$ in the
current dynamics. The Custodian operator $K$ corrects the pathology by
inducing $\kappa(K(B)) \geq 0$.
\end{theorem}

\begin{proof}
\emph{Pathology iff negative curvature.} By Definition~\ref{def:attractor_pathology},
pathology requires $\mu(\Acc_{\mathrm{after}}(B)) < \mu(\Acc_{\mathrm{before}}(B))$,
which is exactly $\kappa(B)<0$.

\emph{$K$ corrects by inducing $\kappa \geq 0$.} From
Theorem~\ref{thm:custodian_curvature}, $K$ breaks the pathology when $K \circ T(B)
\not\subseteq B$, meaning at least one trajectory escapes $B$ per step. The
accessible neighborhood therefore does not shrink: $\kappa(K(B)) \geq 0$. \qed
\end{proof}

\begin{remark}
This reframes the Custodian from a safety mechanism to a geometric necessity.
In any system maintaining accessibility over time, regions of persistently
negative curvature are cognitively pathological: they trap the system,
reducing its future degrees of freedom without recovery. The Custodian is
the runtime operator that maintains non-negative expected curvature on the
accessibility manifold. Its three regimes (allow / throttle / block) correspond
to $\kappa(B) > -\varepsilon$ (normal), $-C \leq \kappa(B) \leq -\varepsilon$
(emerging trap), and $\kappa(B) < -C$ (deep trap requiring full escape).
\end{remark}

\subsection{The \texttt{.m8} Format as File-Level Semantic Projection}

\begin{definition}[Semantic Projection File]
\label{def:semantic_file}
A \emph{semantic projection file} for $x \in X$ under $\Proj: X \to M$ is
a data structure $\mathtt{f}(x)$ such that: (i) $\Proj(x)$ is losslessly
recoverable from $\mathtt{f}(x)$; (ii) for all constraints $c \in
\mathcal{C}_{\mathrm{rel}}$, membership $x \in \Acc(c)$ is reconstructable
from $\mathtt{f}(x)$; and (iii) $|\mathtt{f}(x)|$ is minimized subject to
(i) and (ii).
\end{definition}

\begin{proposition}[\texttt{.m8} is a Semantic Projection File]
\label{prop:m8_semantic}
The \texttt{.m8} format stores the tuple $(id, A, \omega, \phi, I)$ in 32
bytes per memory. Here $id$ recovers $\Proj(x)=m$, and $(A,\omega,\phi,I)$
encodes the fiber coordinate sufficient for behavioral identity
(Proposition~\ref{prop:identity_constraints}). It therefore realizes
Definition~\ref{def:semantic_file}.
\end{proposition}

\begin{proof}
Condition (i): $id$ recovers $m$. Condition (ii): by
Proposition~\ref{prop:wave_dof}, $(A,\omega,\phi,I,D)$ encodes the five
degrees of freedom of the constraint perturbation; by
Proposition~\ref{prop:identity_constraints}, this determines membership in
all relevant admissible regions. Condition (iii): Theorem~\ref{thm:map_territory}
guarantees that $x$ contains information beyond $m$; the wave tuple is the
minimal five-parameter encoding of that additional information. \qed
\end{proof}

\begin{remark}
A conventional file stores content: elements of $M$, the projected space.
An \texttt{.m8} file stores constraint structure: the fiber coordinate, the
un-projected information. The 99\% compression figure is therefore not
primarily a technical achievement but an algebraic consequence:
Theorem~\ref{thm:map_territory} guarantees that $|X| \gg |M|$, so representing
only the five-parameter constraint tuple rather than the full $x \in X$ achieves
compression proportional to the ratio $|X|/|M|$, which is large for high-entropy
memories. The Marqant layer handles residual string-level redundancy of the
projected content.
\end{remark}

\subsection{The Master Derivation}

\begin{theorem}[MEM$|$8 is a Partial Cognitive Runtime]
\label{thm:mem8_runtime}
MEM$|$8 is a partial cognitive runtime for $(\Conf, \Acc, \Proj)$, with the
following theory-to-runtime correspondences:

\begin{center}
\renewcommand{\arraystretch}{1.4}
\begin{tabular}{@{}p{5.5cm}p{5.5cm}@{}}
\toprule
\textbf{Theory primitive} & \textbf{Runtime implementation} \\
\midrule
Constraint field $\Acc$ & Wave grid $(256^2 \times 65536)$ \\
Constraint perturbation  & Memory wave $M_{xyz}(t)$ \\
Overlapping constraint fields & Wave interference (Thm.~\ref{thm:wave_natural}) \\
Constraint relaxation / entropy & Exponential decay $D(t,\tau)$ (Thm.~\ref{thm:exponential_entropy}) \\
Admissibility stack      & Reactive layer hierarchy (Thm.~\ref{thm:reactive_stack}) \\
CLIO inference           & Noise floors, temporal blankets, attention weights \\
Curvature control        & Custodian operator (Thm.~\ref{thm:custodian_curvature}) \\
Semantic projection      & \texttt{.m8} format (Prop.~\ref{prop:m8_semantic}) \\
Chain of Memory          & Wave constraint preservation, not state replay \\
\bottomrule
\end{tabular}
\end{center}
\end{theorem}

\begin{proof}
Each row is established by the cited result. Condition (i) of
Definition~\ref{def:cognitive_runtime}: every constraint is representable as a
wave pattern. Condition (ii): $\Acc(c)$ is computable from the interference
pattern in bounded time. Condition (iii): entropy is estimable from amplitude
distribution (high amplitude = strong constraint = low entropy). Condition
(iv): composition is superposition, relaxation is temporal decay, collapse is
reactive bypass. The runtime is partial: RSVP's global scalar field is not
fully numerically implemented; Spherepop's full containment calculus is
approximated by the urgency-ordered special case; CLIO search is approximated
by threshold filtering. \qed
\end{proof}

\begin{corollary}[Memory = Persistent Modification of Future Accessibility]
\label{cor:memory_theorem}
The equation
\[
  \text{memory} = \text{persistent modification of future accessibility}
\]
is not a definition within the constraint-first framework. It is a theorem: the
unique characterization of cognitively effective memory
(Proposition~\ref{prop:mem_effective}), derived from the accessibility ontology
(Definition~\ref{def:acc_ontology}). When Chenoweth and Chenoweth state this
as a ``final principle''~\cite{mem8_2025}, they are identifying---from an
engineering starting point---the primitive that the present framework derives
from ontological first principles.
\end{corollary}

\begin{corollary}[Consciousness as Recursively Regulated Accessibility]
\label{cor:consciousness_acc}
Consciousness, modeled as global availability of constraint-weighted projections
to deliberative processing, is in MEM$|$8 not the presence of waves alone.
It is the system's capacity to regulate which wave-encoded constraints enter
the active projection manifold:
\[
  \text{consciousness} = \text{recursively regulated accessibility.}
\]
The recursive qualifier reflects that the regulation mechanism is itself subject
to the same constraint-accessibility structure, formalized by the chain of
memory (Definition~\ref{def:chain_of_memory}).
\end{corollary}

\subsection{Critical Assessment}

Theorem~\ref{thm:mem8_runtime} establishes structural correspondence; it does
not validate the empirical claims of~\cite{mem8_2025}. The reported
performance figures (973$\times$ insertion speedup, 292$\times$ retrieval
speedup) lack disclosed experimental protocols and independent replication.
The consciousness metrics are novel and non-standard; their relationship to
peer-reviewed consciousness frameworks remains unestablished. The wave--neural
oscillation connection is analogical: the computational wave grid does not
demonstrably instantiate neural dynamics in any rigorous sense. The ``sensory
free will'' framing conflates control over \emph{access} with phenomenal
consciousness; within the present framework, the architecture establishes that
the system has partial control over which trajectories enter the projection
manifold $M$---agency in the sense of Corollary~\ref{cor:consciousness_acc}
but not phenomenal experience without additional argument.

None of this diminishes the structural convergence. An engineering effort
starting from neural oscillation theory and wave physics independently arrived
at constraint-shaped forgetting, attention as accessibility modulation, reflexive
collapse under narrow admissibility, curvature-controlled attractors, and
identity as constraint rather than state preservation. This convergence
constitutes evidence that the constraint-first architecture reflects something
about the conceptual structure of cognition, not merely an idiosyncratic choice
of ontology.


\subsection{Distributed Mixture-of-Experts as Accessibility Compression}
\label{subsec:dmoe}

The preceding subsections established MEM$|$8 as a partial cognitive runtime
for the constraint-first framework. A recent result in machine learning
provides concrete empirical confirmation that large computational systems
naturally organize around the same structure. The dMoE architecture of
Feng et al.~\cite{dmoe_2026} reduces a diffusion language model's active
expert count from 69.5 to 14.6 while retaining 99.11\% of performance---a
compression ratio of $4.77\times$ with near-zero functional loss. This
section argues that dMoE is not merely an engineering optimization but a
computational instantiation of the Map-Territory Projection Theorem
(Theorem~\ref{thm:map_territory}) and the projection entropy framework
(Section~\ref{sec:entropy}).

\paragraph{The conventional model as process ontology.}
A standard transformer with token-level MoE routing embodies a process
ontology in the sense of Section~\ref{sec:ontological_hierarchy}: computation
is a sequence of token-level operations, and the routing decision is an
individual event attached to each token. The question the model asks at
each step is \emph{which computation should this token perform?}---a local,
trajectory-centric question.

dMoE moves one level down, in exactly the sense of Proposition~\ref{prop:hierarchy}.
Instead of asking which computation each token should perform, it first
constructs a shared routing structure---the block-level expert distribution
$\widetilde{S}_{\mathrm{block}}$---for the entire block, and only then
permits token trajectories inside that structure. The question becomes:
\emph{what is the admissible region of expert space for this block context?}
Token routing is then constrained navigation within that region.

\paragraph{The block distribution as accessibility field.}
The block-level expert distribution is formed by aggregating token-level
router scores:
\[
  S_{\mathrm{block}} = \bigoplus_{i \in B} \hat{s}_i, \qquad
  \hat{s}_i = \mathrm{TopKMask}(s_i, k),
\]
and normalizing to a probability vector $\widetilde{S}_{\mathrm{block}} \in
\Delta^{|E|-1}$. The coreset is then selected by top-$p$ thresholding:
$\mathcal{C} = \mathrm{Top\text{-}P}(\widetilde{S}_{\mathrm{block}}, p)$.

This is an accessibility field in the sense of Definition~\ref{def:acc_ontology}:
it assigns to the current block configuration an admissible set of
computational resources $\mathcal{C}$, and all token-level computation is
constrained to navigate within $\mathcal{C}$. The correspondence with the
framework is exact:

\begin{center}
\renewcommand{\arraystretch}{1.35}
\begin{tabular}{@{}ll@{}}
\toprule
\textbf{Beyond Process} & \textbf{dMoE} \\
\midrule
Constraint field $\mathcal{A}$ & Block-level expert distribution $\widetilde{S}_{\mathrm{block}}$ \\
Admissible trajectories & Expert routes within coreset $\mathcal{C}$ \\
Object (invariant basin) & Stable expert subset under repeated blocks \\
Projection $\pi: X \to M$ & Top-$p$ coreset selection \\
Fiber degeneracy $\pi^{-1}(m)$ & Multiple equivalent token-level routes \\
Projection entropy & Routing redundancy across the full expert space \\
\bottomrule
\end{tabular}
\end{center}

\paragraph{The Map-Territory Theorem in a working model.}
The Map-Territory Projection Theorem (Theorem~\ref{thm:map_territory}) states
that any surjective map $\pi: X \to M$ with $|X| > |M|$ has nontrivial
fibers: distinct elements of $X$ map to the same element of $M$, and the
information distinguishing them is invisible in $M$. dMoE's core empirical
result is a direct computational demonstration of this theorem.

Let $X$ be the space of all possible token-level routing assignments across
a block (combinatorially large, scaling as $\binom{|E|}{k}^{|B|}$) and
$M$ be the space of block-level routing assignments constrained to the
coreset (much smaller, scaling as $\binom{|\mathcal{C}|}{k}^{|B|}$ with
$|\mathcal{C}| \approx 14$ versus $|E| \approx 70$). The projection
$\pi: X \to M$ maps full routing assignments to their coreset-constrained
counterparts.

The nontrivial fibers are the distinct full-routing assignments that produce
the same coreset-constrained output: they are \emph{functionally equivalent}
under the projection, because the model cannot distinguish between them at
the block level. The performance result (99.11\% retention) confirms that
these fibers are not only large but \emph{functionally degenerate}: the
information that distinguishes elements within a fiber carries almost no
functional significance. Discarding the fiber coordinate---collapsing to
the coreset manifold---loses less than 1\% of performance.

\begin{proposition}[dMoE Instantiates the Projection Theorem]
\label{prop:dmoe_projection}
The dMoE result is a computational instantiation of Theorem~\ref{thm:map_territory}:
the full routing space $X$ projects onto the coreset space $M$ via
$\pi = \mathrm{Top\text{-}P} \circ \mathrm{Normalize} \circ \bigoplus$,
and the fibers $\pi^{-1}(m)$ are large but functionally degenerate, as
confirmed by near-zero performance loss under coreset restriction.
\end{proposition}

\begin{proof}
By Theorem~\ref{thm:map_territory}, $|X| > |M|$ guarantees nontrivial fibers.
The empirical result establishes that the fibers are \emph{functionally}
nontrivial in size (many token routes map to the same coreset) but
\emph{informationally} degenerate (they are functionally equivalent for the
model's task). This is exactly the projection theorem's content: the
information distinguishing fiber elements---which specific routes tokens
took within the coreset---is invisible in the observable output. \qed
\end{proof}

\paragraph{Memory as coreset storage, not route storage.}
The connection to MEM$|$8 is now precise. Theorem~\ref{thm:mem8_runtime}
established that MEM$|$8 stores constraint structure (the fiber coordinate
sufficient for behavioral identity) rather than full state trajectories.
dMoE suggests what this means concretely in a large language model context:
effective memory stores the block accessibility field $(S_{\mathrm{block}},
\mathcal{C})$ rather than the individual routing decisions $\{R_i\}_{i \in B}$.
Future inference reconstructs the admissibility manifold from the stored
coreset, not from replaying stored routes.

This is the computational realization of the Chain of Memory principle
(Definition~\ref{def:chain_of_memory}): identity persists through
constraint-preserving reconstruction, not state preservation. The stored
coreset $\mathcal{C}$ carries the constraint structure---which experts are
admissible---while the specific routes within $\mathcal{C}$ are fiber
elements: interchangeable, degenerate, and safely discarded.

\begin{corollary}[Effective Memory = Admissibility Structure, Not Route History]
\label{cor:mem_admissibility}
A system that stores the block-level accessibility field
$(S_{\mathrm{block}}, \mathcal{C})$ rather than individual token routes
$\{R_i\}_{i \in B}$ retains the functionally relevant constraint structure
while discarding fiber-degenerate routing details. The dMoE result confirms
empirically that this compression is achievable with $4.77\times$ reduction
in stored routing complexity and near-zero functional loss.
\end{corollary}

\paragraph{Top-$p$ as CLIO in expert space.}
The top-$p$ coreset selection is formally isomorphic to the CLIO sparse
admissibility operator (Theorem~\ref{thm:admissibility_coherence}).
CLIO restricts inference to admissible queries given a document context.
dMoE restricts routing to admissible experts given a block context.
Both are instances of the same operation---sparse projection onto a
tractable admissible subset---applied to different substrates: inferential
state space in CLIO, expert index space in dMoE.

The threshold $p$ in dMoE corresponds to the coherence threshold $\theta$
in Chain of Memory: it determines how much of the accessibility structure
must be preserved for the projection to be considered semantically coherent.
A lower $p$ means more aggressive compression; a lower $\theta$ means
looser identity preservation. Both trade fiber detail for computational
tractability.


%% ══════════════════════════════════════════════════════════════════════════════
\section{The Implementation Tower}
\label{sec:implementation_hierarchy}
%% ══════════════════════════════════════════════════════════════════════════════

The full derivation assembles into a layered tower in which each level is the
unique realization of the level above under tractability and continuity
constraints.

\begin{definition}[The Constraint-First Implementation Tower]
\label{def:impl_tower}
\begin{center}
\renewcommand{\arraystretch}{1.5}
\begin{tabular}{@{}cll@{}}
\toprule
\textbf{Level} & \textbf{Framework} & \textbf{Domain} \\
\midrule
L1 & Semantic Infrastructure & Ontology: constraints, accessibility, projection \\
L2 & CLIO           & Inference: sparse admissibility search \\
L3 & Spherepop      & Evaluation: urgency-ordered containment collapse \\
L4 & Chain of Memory & Continuity: constraint-preserving reconstruction \\
L5 & RSVP           & Field dynamics: scalar/vector accessibility evolution \\
L6 & MEM$|$8        & Runtime: wave grids, decay, reactive layers, compression \\
\bottomrule
\end{tabular}
\end{center}
Each level is a projection of L1 onto a specific computational or physical screen.
\end{definition}

\begin{theorem}[Each Level is Uniquely Derived from the Level Above]
\label{thm:tower_derivability}
Each level L$(n+1)$ is the unique (up to implementation choice) realization of
L$(n)$'s abstract principles subject to the constraints of tractability,
continuity, and minimality.
\end{theorem}

\begin{proof}
\emph{L1 $\to$ L2.} Given that cognition is accessibility navigation, inference
must restrict to admissible regions (exhaustive search over $X$ is intractable
by cardinality). CLIO is the minimal inference principle preserving semantic
coherence (Theorem~\ref{thm:admissibility_coherence}) while bounding search.

\emph{L2 $\to$ L3.} Given CLIO-style restricted inference, evaluation requires
a total order. Urgency-ordered collapse is the unique latency-optimal order:
when $|\Acc(c)|=1$, no other order can improve the output
(Proposition~\ref{prop:reactive_spherepop}).

\emph{L3 $\to$ L4.} Given ordered collapse, persistence through collapse
requires constraint-based identity (states are destroyed by collapse;
Proposition~\ref{prop:identity_constraints}). Chain of Memory is the minimal
continuity structure satisfying this.

\emph{L4 $\to$ L5.} Given constraint-preserving reconstruction, a field
language is needed for how the accessibility landscape evolves between collapse
events. RSVP is the unique Lorentz-compatible scalar-vector formalism for
accessibility evolution (Proposition~\ref{prop:rsvp_eom}).

\emph{L5 $\to$ L6.} Given a field language, a runtime must encode constraint
perturbations numerically. Theorem~\ref{thm:wave_natural} establishes that
complex wave representation is the unique natural encoding for overlapping
constraint fields. Theorem~\ref{thm:exponential_entropy} establishes exponential
decay as the entropy-optimal relaxation schedule. Together they uniquely
determine the wave grid architecture. Reactive layers, Custodian, and
\texttt{.m8} follow from Theorems~\ref{thm:reactive_stack},
\ref{thm:custodian_curvature}, and Proposition~\ref{prop:m8_semantic}. \qed
\end{proof}

\begin{remark}[The Tower as a Processor--Language Stack]
The six-level relationship is precisely analogous to a compiler/processor stack:

\begin{center}
\renewcommand{\arraystretch}{1.4}
\begin{tabular}{@{}ll@{}}
\toprule
\textbf{CS analogy} & \textbf{Constraint-first tower} \\
\midrule
Specification language   & Semantic Infrastructure (L1) \\
Type system / inference  & CLIO (L2) \\
Evaluation strategy      & Spherepop (L3) \\
Memory model             & Chain of Memory (L4) \\
Instruction set / ABI    & RSVP (L5) \\
Processor microarchitecture & MEM$|$8 (L6) \\
\bottomrule
\end{tabular}
\end{center}

Just as a program written in a high-level language is ultimately executed by
transistor-level operations, a cognitive process described at the level of
accessibility ontology is ultimately executed by wave-grid operations. The
levels do not contradict each other; each is a compiled form of the one
above. This is what it means to say that MEM$|$8 is a \emph{partial machine
realization} of the constraint-first framework---not a compatible or analogous
system, but an execution model for its primitives.
\end{remark}

The claim in the preface of \emph{Semantic Infrastructure}---that the frameworks
are ``five projections of one theory onto five different observable screens''---is
therefore a theorem. Theorem~\ref{thm:tower_derivability} establishes that the
projections are deductive, not analogical: each screen is the unique
realization of the theory's primitives at that computational level.


%% ══════════════════════════════════════════════════════════════════════════════
\section{Conclusion: Everything Navigates}
\label{sec:conclusion}
%% ══════════════════════════════════════════════════════════════════════════════

The arguments of the preceding sections can now be gathered into a single
proposition.

\begin{theorem}[Master Theorem]
\label{thm:master}
Let any system be described by a triple $(\Conf, \Acc, \Proj)$ where $\Conf$
is a configuration space, $\Acc: \Conf \to \mathcal{P}(\Conf)$ a constraint
field, and $\Proj: \Conf \to M$ a projection to an observable space $M$. Then:
\begin{enumerate}[label=(\roman*)]
  \item All stable structures are basins of $\Acc$
    (Proposition~\ref{prop:hierarchy}).
  \item All processes are trajectories through $\Acc$
    (Definition~\ref{def:acc_ontology}).
  \item All information loss is projection entropy
    (Proposition~\ref{prop:boltzmann}).
  \item All coordination is stigmergic modification of $\Acc$
    (Proposition~\ref{prop:stigmergy_unification}).
  \item All inference is admissibility-constrained trajectory selection
    (Theorem~\ref{thm:admissibility_coherence}).
\end{enumerate}
Therefore, the appropriate governing slogan is not Heraclitus's
\emph{``everything flows''} but: \emph{everything navigates a changing
landscape of admissible possibilities.}
\end{theorem}

\begin{proof}
Each clause follows from the corresponding theorem or proposition cited.
Together, they constitute an exhaustive account of structure, process,
information, coordination, and inference within the framework. The slogan
collects the common pattern: in each case, the relevant phenomenon is not
the trajectory taken but the accessibility landscape within which trajectory
selection occurs.
\qed
\end{proof}

\begin{corollary}[Unity of the Frameworks]
Spherepop, RSVP, CLIO, Admissibility Lab, Yarncrawler, Semantic
Infrastructure, and MEM$|$8 are projections of a single underlying theory---a
geometric theory of constrained accessibility---onto seven different observable
screens: computation, cosmology, inference, document structure, narrative
coherence, civilizational organization, and artificial cognitive runtime.
\end{corollary}

\begin{proof}
Each framework instantiates the triple $(\Conf, \Acc, \Proj)$: Spherepop with
$\Conf = $ containment structures, $\Acc = $ admissible Pop operations,
$\Proj = $ string linearization; RSVP with $\Conf = $ field configurations,
$\Acc = $ accessibility gradients, $\Proj = $ cosmological observables; CLIO
with $\Conf = $ document states, $\Acc = $ admissible queries, $\Proj = $
answer space; Admissibility Lab with $\Conf = $ semantic neighborhoods,
$\Acc = $ perturbation admissibility, $\Proj = $ embedding clusters;
Yarncrawler with $\Conf = $ narrative trajectory space, $\Acc = $ coherent
continuations, $\Proj = $ surface text; Semantic Infrastructure with
$\Conf = $ full knowledge-state space, $\Acc = $ admissible inferential moves,
$\Proj = $ observable propositional content. The structural isomorphism of all
six instantiations confirms a common origin.
\qed
\end{proof}

The history of philosophy has debated being versus becoming for two and a half
millennia. The present proposal does not resolve that debate by choosing a
side. It relocates the question. Before there can be stable beings or flowing
becomings, there must be a geometry---an accessibility landscape---that makes
some states persistent and some transitions possible. That geometry is the
ground of both. It is not a thing, not an event, not a flow. It is the
structure of what can happen from here.

%% ── Bibliography note ─────────────────────────────────────────────────────────
\bigskip
\noindent\rule{\textwidth}{0.4pt}
\smallskip
\noindent{\small
Principal sources for the process philosophy sections: A.\,N.\ Whitehead,
\emph{Process and Reality} (1929); Nicholas Rescher, \emph{Process Philosophy}
(1996); Heraclitus, fragments (DK 22). For information geometry and constraint
ontology: S.-i.\ Amari, \emph{Information Geometry and Its Applications}
(2016); J.\ Ladyman and D.\ Ross, \emph{Every Thing Must Go} (2007). For
stigmergy: P.-P.\ Grass\'{e} (1959); E.\ Bonabeau et al., \emph{Swarm
Intelligence} (1999). For MEM$|$8: C.\,M.\ Chenoweth and A.\,A.\ Chenoweth,
``MEM$|$8: A Wave-Based Cognitive Architecture for Multimodal Memory Integration
and Consciousness Simulation,'' Zenodo preprint, DOI:~10.5281/zenodo.16436298
(2025). All formal frameworks (RSVP, Spherepop, CLIO, Admissibility Lab,
Yarncrawler, Semantic Infrastructure) are the author's own; see associated
monographs at \texttt{github.com/standardgalactic}.
}


\newpage
%% ── Bibliography ────────────────────────────────────────────────────────────
\begin{thebibliography}{9}

\bibitem{whitehead1929}
A.\,N.\ Whitehead, \emph{Process and Reality: An Essay in Cosmology}.
Macmillan, New York, 1929.

\bibitem{rescher1996}
N.\ Rescher, \emph{Process Metaphysics: An Introduction to Process Philosophy}.
SUNY Press, Albany, 1996.

\bibitem{amari2016}
S.\ Amari, \emph{Information Geometry and Its Applications}.
Springer, Tokyo, 2016.

\bibitem{ladyman2007}
J.\ Ladyman and D.\ Ross, \emph{Every Thing Must Go: Metaphysics Naturalised}.
Oxford University Press, Oxford, 2007.

\bibitem{grasse1959}
P.-P.\ Grass\'{e}, ``La reconstruction du nid et les coordinations
inter-individuelles chez \emph{Bellicositermes natalensis} et \emph{Cubitermes}
sp.\ La th\'{e}orie de la stigmergie,'' \emph{Insectes Sociaux}, 6:41--83, 1959.

\bibitem{bonabeau1999}
E.\ Bonabeau, M.\ Dorigo, and G.\ Theraulaz, \emph{Swarm Intelligence: From
Natural to Artificial Systems}. Oxford University Press, New York, 1999.

\bibitem{baars1988}
B.\,J.\ Baars, \emph{A Cognitive Theory of Consciousness}.
Cambridge University Press, Cambridge, 1988.

\bibitem{franklin2014}
S.\ Franklin et al., ``LIDA: A systems-level architecture for cognition,
emotion, and learning,'' \emph{IEEE Transactions on Autonomous Mental
Development}, 6(1):19--41, 2014.

\bibitem{dmoe_2026}
S.\ Feng, Z.\ Chen, G.\ Fang, X.\ Ma, and X.\ Wang, ``dMoE: dLLMs with
Learnable Block Experts,'' arXiv:2605.30876, May 2026.

\bibitem{mem8_2025}
C.\,M.\ Chenoweth and A.\,A.\ Chenoweth, ``MEM$|$8: A Wave-Based Cognitive
Architecture for Multimodal Memory Integration and Consciousness Simulation,''
Zenodo preprint, DOI:~10.5281/zenodo.16436298, July 2025.

\end{thebibliography}

\end{document}

%%% ─── REMOVE OLD END AND INJECT NEW SECTIONS ───
